% page 411 du fac-simile (image p-410 du scan)
% bloc 162.6 x 246.3 mm ; marge gauche 39.4 mm
\pgc{17.780mm}{0.000mm}{
\l{\cel{54}403}
\l{}
\l{\sou{odontoblasto}. (\sou{biol.)}\cel{2}Grosa celulo, kun nukleo tre voluminoza,}
\l{\cel{3}quan onu renkontras en la parto periferia di la pulpo den-}
\l{\cel{3}tala, e di qua la rolo specala esas sekrecar ivoro. -}
\l{}
\l{\sou{odorar}.\cel{2}(\sou{netrans}.)(\sou{aludante ula korpi}). Exhalar ulo qua im-}
\l{\cel{3}presas partikulare la olfakt-organo. - EFIS.}
\l{}
\l{\sou{ofensar}.\cel{2}(\sou{trans.})\cel{2}Vundar (ulu) pri lua digneso, respektind-}
\l{\cel{3}eso. - EFIS.}
\l{}
\l{\sou{ofertorio}. (\sou{liturgio katol.})\cel{2}I.Deo-ofro di pano e vino en}
\l{\cel{3}la meso-sakrifiko. - II. Prego, muziko, qui akompanas ta}
\l{\cel{3}ofro. - DEFIS.}
\l{}
\l{\sou{ofico}. Funciono administrala o direktala, sive en stato, sive}
\l{\cel{3}en societo o granda entraprezajo. - EFIS.}
\l{}
\l{\sou{oficiar}. (\sou{netrans}.) (\sou{eklezio katol.)}\cel{3}Celebrar la deo-servo}
\l{\cel{3}kun la ceremonii qui devas akompanar lu. -\cel{2}EFIS.}
\l{}
\l{\sou{oficiro}. (\sou{armeo)}.\cel{2}La persono qua havas ofico komandala en}
\l{\cel{3}armeo (terala, marala, aerala). - DEFIRS.}
\l{}
\l{\sou{ofikleido}. (\sou{muziko}). Sufl-instrumento kupra, kun boko-peco,}
\l{\cel{3}qua ne esas altro kam la serpento-korno, quan onu uzis en}
\l{\cel{3}la kirki, a qua onu adjuntis klavi. - DEFIS.}
\l{}
\l{\sou{ofito}. (\sou{mineral.}) Marmoro obskur-verda, kun strii flava, quan}
\l{\cel{3}onu renkontras en Pirenei. - DEF.}
\l{}
\l{ofrar. (\sou{trans., ad)}\cel{2}Igar (ulo, o, metaf., ulu) disponebla da}
\l{\cel{3}altru, sen ke ica demandis lo. - DEFIS.}
\l{}
\l{\sou{ofte.}\cel{1}Qua eventas granda-quanta-foye, ma ne ye intertempi}
\l{\cel{3}inter-egala - nur intermite. - DE.}
\l{}
\l{oftalmio. (\sou{patol.)}\cel{2}Morbo inflamurala di la okulo. - DEFIS.}
\l{}
\l{oftalmologio.\cel{2}Cienco qua traktas la strukturo e la funciono}
\l{\cel{3}di la okuli e la morbi qui afektas li. -}
\l{}
\l{ogivo. I. (\sou{arkeol.)}\cel{3}Arko ruptita quan formacas, en la arki-}
\l{\cel{3}tekturo gotika, la nervaturi di la travei di la vulto, inter-}
\l{\cel{3}krucumante diagonale ye la kolmo. - II. (\sou{arkitekt.)}\cel{1}Arkado}
\l{\cel{3}ek du arki\cel{1}\sou{qui}\cel{1}interkrucumas formacante angulo kurvo-linea,}
\l{\cel{3}akuta. - DEFIS.}
\l{}
\l{oglar. (trans.) Regardar furtatre, quaze per la komisuri di}
\l{\cel{4}singla okulo (ulo) quan onu deziregas. - E.}
\l{}
\l{\cel{1}\sou{\textquotedbl{}ohm\textquotedbl{}}\cel{2}(\sou{elektr.}) Unajo praktikala di la rezisto elektrala =}
\l{\cel{4}l09 unaji elektromagnetala C.G.S.\cel{3}DEFIRS.}
\l{}
\l{\cel{1}- oid.\cel{2}Sufixo ciencala : \textquotedbl{}qua havas la sama formo kam\textquotedbl{} (o}
\l{\cel{4}formo analoga).}
}
